\chapter{Methodology}
\label{chapter:methodology}

This chapter describes how the theoretical principles established in
Chapter~\ref{chapter:theoretical_background} are realised in a working ascent-trajectory
simulator, and how that simulator is driven by the optimisation back-ends used
throughout this work. The presentation is kept at the level of architecture,
data flow, and configuration: it explains \emph{where} and \emph{how} each model
is implemented without reproducing source code. The chapter opens with the
software architecture, then follows the natural order of a single evaluation ---
mission phases, the integrated state and numerical scheme, the reference-frame
strategy, the equations of motion, and the environmental and force models. It
then catalogues the nine guidance laws and the four optimisation architectures
available, and closes with the configuration system, the baseline vehicle and
mission, and a note on running the software.

% =============================================================================
\section{Software Architecture Overview}
\label{sec:architecture}

The code is organised as a layered dependency graph. At the bottom sit leaf data
and models that depend on nothing else: the physical constants and vehicle
specifications, and the environment models (gravity, atmosphere, Earth rotation).
Above them the configuration layer exposes every user-facing switch. The guidance
laws are also leaves --- pure functions that map a state to a steering command ---
so that dependencies always run from the simulation core towards the guidance
modules and never the reverse, which keeps each law testable in isolation. At the
centre is the physics hub \texttt{rocket\_ascent.py}, which imports all guidance
and auxiliary modules and integrates the equations of motion; the solvers and
optimisers wrap the hub, and the plotting layer consumes plain result arrays,
decoupled from the solvers.

\begin{figure}[h!]
\centering
\begin{minipage}{0.92\linewidth}
\small
\begin{verbatim}
main.py
  |-- Simulation/solver.py --------------------+
  |-- Simulation/pso_coast_solver.py           |  all import ->
  |-- Simulation/direct_pso_solver.py          |  Simulation/rocket_ascent.py
  |     \-- imports pso_coast_solver            |  (the hub: imports ALL
  |-- Simulation/indirect_pso_solver.py -------+   Guidance/* + Auxiliary/*)
  |         \-- imports Guidance/indirect_pmp_guidance.py
  |-- Auxiliary/{constants, gravity, atmosphere, earth_rotation, rocket_specs}.py
  |-- Input_File/simulation_parameters.py   (imported almost everywhere)
  \-- Plots/new_plot_runner.py -> Plots/new_metrics/*   (decoupled)
\end{verbatim}
\end{minipage}
\caption{Surface-level module dependency tree. \texttt{rocket\_ascent.py} is the
central hub; the guidance modules are leaves, so dependencies run
Simulation~$\rightarrow$~Guidance and never the reverse.}
\label{fig:dependency_tree}
\end{figure}

The execution flow is an outer--inner loop. An outer optimiser proposes a
candidate design vector; the inner simulation engine integrates the trajectory,
evaluating at every step the environment models, the active guidance law, and the
equations of motion, with event detection triggering phase transitions and
cut-off; the resulting trajectory is scored by an objective function and returned
to the optimiser (Figure~\ref{fig:flowchart}). The same inner engine serves all
four optimisation architectures of Section~\ref{sec:optimization}; only the design
vector, the termination event, and the objective differ.

% NOTE (author action): FlowChart.png should be redrawn to show the four
% optimisation back-ends branching on COAST_METHOD / GUIDANCE_MODE, replacing the
% single brute-force loop of the original figure.
\begin{figure}[h!]
    \centering
    \includegraphics[width=1\linewidth]{Figures/FlowChart.png}
    \caption{Flight-dynamics simulation and optimisation flowchart: the outer
    optimiser feeds a design vector to the inner engine, which returns a scalar
    cost. The inner engine is guidance- and architecture-agnostic.}
    \label{fig:flowchart}
\end{figure}

% =============================================================================
\section{Mission Decomposition into Phases}
\label{sec:phases}

A single trajectory evaluation is resolved as an ordered sequence of mission
phases, each propagated by a separate call to the numerical integrator and
chained by event detection:

\begin{enumerate}
  \item \textbf{Stage-1 powered ascent.} From lift-off until first-stage
        propellant depletion (Main Engine Cut-Off, MECO): an initial vertical
        rise, the pitchover kick manoeuvre, and atmospheric flight under
        aerodynamic drag. This phase is produced by
        \texttt{rocket\_ascent.run\_stage1}.

  \item \textbf{Stage separation and coast.} An instantaneous drop of the
        first-stage inert mass followed by a short ballistic coast to
        second-stage ignition.

  \item \textbf{Stage-2 powered ascent with guidance.} The active guidance law
        (one of the nine modes of Section~\ref{sec:guidance}) shapes the
        trajectory until the architecture-dependent cut-off.

  \item \textbf{Coast and orbit insertion.} Depending on the mission
        architecture (Section~\ref{sec:optimization}), the vehicle either coasts
        to apogee for an impulsive circularisation burn, or is inserted directly
        when the powered arc reaches circular-velocity conditions.
\end{enumerate}

The thrust--coast--thrust arc structure of the PSO architectures is assembled
inside the solvers, which reuse the same right-hand side as
\texttt{rocket\_ascent.run} with the thrust set to zero on coast arcs.

% =============================================================================
\section{State Vector and Numerical Integration}
\label{sec:state_integration}

\subsection{State Vector}
\label{ssec:state_vector}

The base state integrated by the solver is
\begin{equation}
  \mathbf{x} = \bigl[\,s,\; r,\; v,\; \gamma,\; m\,\bigr]^{T},
  \label{eq:state_base}
\end{equation}
with downrange distance $s$, geocentric radius $r$, rotating-frame speed $v$,
flight-path angle $\gamma$, and mass $m$. When Earth-rotation effects are enabled
the geocentric latitude $\phi$ is appended,
\begin{equation}
  \mathbf{x} = \bigl[\,s,\; r,\; v,\; \gamma,\; m,\; \phi\,\bigr]^{T},
  \label{eq:state_rot}
\end{equation}
so that the pseudo-force accelerations can be evaluated at the current latitude.
The local heading (azimuth) $\psi$ is \emph{not} propagated as a dynamical state:
it is held constant at the launch azimuth throughout the ascent, since the planar
model carries no cross-range degree of freedom. A dynamic heading state was used
in an earlier version of the model and is retained only as a deprecated option;
reinstating full three-dimensional heading tracking is left as future work.

\subsection{Numerical Integration}
\label{ssec:integration_sim}

Each phase is propagated with the explicit Runge--Kutta~4(5) Dormand--Prince
scheme provided by \texttt{scipy.integrate.solve\_ivp}~\cite{dormand1980family,
virtanen2020scipy}, using adaptive step-size control, a tight tolerance, a bounded
maximum step, and dense output for post-processing. Phase transitions are located
by \emph{event functions} --- scalar functions $g(t,\mathbf{x})$ whose sign change
marks a transition. When the integrator brackets a zero-crossing, the exact
crossing time is found by root-finding (Brent's method) and the phase is
terminated~\cite{brent1973algorithms}. Events include propellant depletion
(MECO), stage separation, atmosphere exit, target-apogee matching (SECO), the
circular-velocity cut-off used by the direct architecture, and ground-collision
safeguards. Chaining phases through events lets the mission timeline emerge from
the physics rather than being prescribed in advance.

% =============================================================================
\section{Reference-Frame Strategy}
\label{sec:sim_frames}

The simulator uses the mixed-frame strategy motivated in
Section~\ref{sec:reference_frames}: ascent dynamics quantities ($v$, $\gamma$,
and the heading decomposition) are expressed in a rotating (ECEF-like) frame,
which makes the treatment of atmosphere-relative drag natural; orbital energy and
orbital elements are evaluated in an Earth-Centred Inertial (ECI) frame, which is
required for consistent event detection (target-apogee matching) and end-of-run
orbit characterisation; and the Coriolis and centrifugal pseudo-accelerations are
first formed in the local East--North--Up basis before being projected onto the
flight-dynamics directions.

During powered ascent and early coasting the integration proceeds in the rotating
frame with optional pseudo-force corrections. Once the vehicle reaches the final
ballistic coast (post-SECO), the propagation switches to pure inertial two-body
dynamics, and the pseudo-force terms are deactivated so that no spurious
accelerations are introduced when the motion is governed solely by gravity.

Whenever orbital elements are required, the rotating-frame velocity is converted
to its inertial counterpart with the component relations derived in
Eqs.~\eqref{eq:ecef_eci}: only the eastward component gains the surface rotation
speed $\omega_E r\cos\phi$. The launch azimuth is computed from the target
inclination and launch latitude using the rotating-Earth correction of
Eq.~\eqref{eq:ground_azimuth}; the simulator applies the corrected ground-relative
azimuth $A_G$ by default and exposes the geometric inertial value $A_I$ as an
option. At end of run the achieved inclination and its drift
$\Delta i = i_{\mathrm{achieved}} - i_{\mathrm{target}}$ follow from
Eq.~\eqref{eq:achieved_inclination}.

% =============================================================================
\section{Implemented Equations of Motion}
\label{sec:eom}

\subsection{Base Dynamics}
\label{ssec:base_eom}

The core system integrated by the simulator, realising the planar model of
Eq.~\eqref{eq:planar_eom}, is
\begin{align}
  \dot s      &= \frac{R_E}{r}\,v\cos\gamma,
               \label{eq:sim_sdot}\\[3pt]
  \dot r      &= v\sin\gamma,
               \label{eq:sim_rdot}\\[3pt]
  \dot v      &= \frac{F_T}{m}\cos\alpha
                 - \frac{F_D}{m}
                 - g(r)\sin\gamma,
               \label{eq:sim_vdot}\\[3pt]
  \dot\gamma  &= \frac{1}{v}\!\left(
                   \frac{F_T}{m}\sin\alpha
                   + \frac{F_L}{m}
                   - \Bigl[g(r) - \frac{v^2}{r}\Bigr]\cos\gamma
                 \right),
               \label{eq:sim_gammadot}\\[3pt]
  \dot m      &= -\frac{F_T}{I_{sp}\,g_0},
               \label{eq:sim_mdot}
\end{align}
where $\alpha$ is the angle of attack (the steering command from the active
guidance law), $F_T$ and $I_{sp}$ are the current thrust and specific impulse,
$F_D$ is aerodynamic drag, $F_L$ is lift (set to zero in every analysis here),
and $g(r)=\mu/r^{2}$. The centripetal term $v^{2}/r$ in
Eq.~\eqref{eq:sim_gammadot} raises the flight-path angle once it exceeds
$g\cos\gamma$. Two safeguards are applied: a small-velocity guard forces
$\dot\gamma=0$ when $v<10^{-6}$~m/s to avoid the $1/v$ singularity, and during
the initial vertical rise (before pitchover) $\dot\gamma$ is held at zero
regardless of the geometric terms.

\subsection{Rotating-Frame Augmentation}
\label{ssec:rot_augmentation}

When Earth rotation is enabled, the Coriolis and centrifugal accelerations
of Eqs.~\eqref{eq:coriolis_enu}--\eqref{eq:centrifugal_enu} are summed in each
ENU direction and projected onto the along-heading and vertical directions
(Eq.~\eqref{eq:pseudo_projection}). The resulting additive corrections to the
velocity and flight-path-angle rates are
\begin{align}
  \Delta\dot v     &= a_{h,\parallel}\cos\gamma + a_{\mathrm{up}}\sin\gamma,
                    \label{eq:sim_dv_pseudo}\\[3pt]
  \Delta\dot\gamma &= \frac{-a_{h,\parallel}\sin\gamma
                      + a_{\mathrm{up}}\cos\gamma}{v},
                    \label{eq:sim_dgamma_pseudo}
\end{align}
so that the integrated rates become $\dot v + \Delta\dot v$ and
$\dot\gamma + \Delta\dot\gamma$; with rotation disabled both corrections vanish
and Eqs.~\eqref{eq:sim_vdot}--\eqref{eq:sim_gammadot} are recovered exactly.

\subsection{Latitude Reconstruction}
\label{ssec:latitude}

Rather than integrating a latitude rate directly, the simulator reconstructs
latitude from the downrange angle by great-circle geometry, which avoids
numerical drift~\cite{BateMuellerWhite1971Fundamentals}:
\begin{equation}
  \sigma = \frac{s}{R_E}, \qquad
  \sin\phi = \sin\phi_0\cos\sigma + \cos\phi_0\sin\sigma\cos\beta_0,
  \label{eq:lat_recon}
\end{equation}
where $\phi_0$ is the launch latitude and $\beta_0$ the inertial launch azimuth.
The latitude rate required by the pseudo-force evaluation follows analytically,
\begin{equation}
  \dot\phi = \frac{-\sin\phi_0\sin\sigma
    + \cos\phi_0\cos\sigma\cos\beta_0}{\cos\phi}\;\frac{\dot s}{R_E}.
  \label{eq:phidot}
\end{equation}

% =============================================================================
\section{Environmental and Force Models}
\label{sec:models}

\subsection{Gravity and Atmosphere}
\label{ssec:gravity_atm}

Gravity is the central inverse-square field $g(r)=\mu/r^{2}$ with
$\mu=3.986\times10^{14}~\mathrm{m^{3}\,s^{-2}}$; the Earth is treated as
spherical, so oblateness ($J_2$) is not modelled. Atmospheric density follows the
exponential profile of Eq.~\eqref{eq:exp_atmosphere}, with sea-level density
$\rho_0=1.225~\mathrm{kg\,m^{-3}}$ and scale height $H=8500$~m, and the dynamic
pressure is $q=\tfrac12\rho v^{2}$. Two criteria are available for declaring
atmosphere exit and activating exo-atmospheric guidance: an altitude threshold
($h>65$~km~\cite{kluever2018space}) or a dynamic-pressure threshold
($q<1000$~Pa, configurable). The selected criterion is a simulation parameter and
is shared with the fairing-jettison logic.

\subsection{Aerodynamic and Propulsive Forces}
\label{ssec:aero_prop}

Aerodynamic drag is evaluated at every step from the dynamic pressure and a
constant drag coefficient $C_D=0.3$ over the reference area
$A=10.52~\mathrm{m^{2}}$~\cite{NASA_Falcon9v12_Datasheet_2018}, using
Eqs.~\eqref{eq:cd} and~\eqref{eq:q_dyn_press}. A lift model is present in the
formulation (the $F_L$ term of Eq.~\eqref{eq:sim_gammadot}) but is neglected in
all analyses ($F_L=0$).

Thrust magnitude $F_T$ and specific impulse $I_{sp}$ are stage-dependent
constants that switch at discrete events (MECO, separation, second-stage
ignition, cut-off). For the first stage, the effective sea-level-to-vacuum
variation is captured by selectable engine modes (\texttt{ISP\_1\_MODE} and
\texttt{THRUST\_1\_MODE}: sea-level, vacuum, average, or altitude-linear). Thrust
acts at the angle of attack $\alpha$ relative to the velocity vector; $\alpha$ is
the primary control and is set at each step by the active guidance law
(Section~\ref{sec:guidance}). Mass depletion obeys Eq.~\eqref{eq:sim_mdot} with
$g_0=9.80665~\mathrm{m\,s^{-2}}$. A no-atmosphere mode (\texttt{INCLUDE\_DRAG}
disabled) removes drag, drops the fairing at lift-off, and forces the
altitude-based exit criterion, giving a clean reference free of drag loss.

% =============================================================================
\section{Guidance Laws Implemented}
\label{sec:guidance}

Nine guidance modes are available, selected by \texttt{GUIDANCE\_MODE}. All share
a common initial pitchover manoeuvre and differ in how the angle of attack
$\alpha$ is commanded thereafter. The scalar-$t_{go}$ modes (Apollo, the tangent
laws, and CPR) share a \texttt{TGO\_ESTIMATOR} option that selects between a
rocket-equation and a PEG-based time-to-go estimate. The atmospheric-arc theory
of each law is developed in Chapter~\ref{chapter:theoretical_background}; this section
documents the implemented form.

\subsection{Common Initial Kick Phase}
\label{ssec:kick_phase}

Regardless of mode, the vehicle rises vertically until time $T_k$ (nominally
$7.5$~s) and then executes a symmetric triangular pitchover over a duration
$\Delta T_k$ (nominally $45$~s):
\begin{equation}
  \alpha(t) =
  \begin{cases}
    0, & t < T_k, \\[2pt]
    \displaystyle\alpha_k\;\frac{2(t - T_k)}{\Delta T_k},
      & T_k \le t < T_k + \tfrac{\Delta T_k}{2}, \\[6pt]
    \displaystyle 2\, \alpha_k\!\left(1 - \frac{t - T_k}{\Delta T_k}\right),
      & T_k + \tfrac{\Delta T_k}{2} \le t < T_k + \Delta T_k, \\[6pt]
    0, & t \ge T_k + \Delta T_k,
  \end{cases}
  \label{eq:kick_profile}
\end{equation}
where $\alpha_k$ is the kick angle. After the kick, and until the atmosphere-exit
criterion is met, $\alpha$ is held at zero; active guidance then engages. The
kick angle (or, equivalently, the pitch-manoeuvre parameter $\gamma_p$) is the
design variable optimised by every architecture of Section~\ref{sec:optimization}.

\subsection{Gravity Turn}
\label{ssec:gt}

The gravity-turn mode applies no active steering after the kick,
$\alpha(t)=0$ for all $t>T_k+\Delta T_k$, so the trajectory is shaped entirely by
gravity and the initial conditions. It is the simplest baseline.

\subsection{Constant Pitch Rate (CPR)}
\label{ssec:cpr}

CPR ramps the commanded pitch angle linearly from its value at guidance start
towards the horizontal at a constant rate $\dot\theta$ fixed once, at engagement,
from the estimated time-to-go~\cite{Ulrich}:
\begin{equation}
  \theta_{\mathrm{cmd}}(t) = \max\!\bigl(\theta_{\mathrm{init}}
        - \dot\theta\,(t - t_{\mathrm{start}}),\; 0\bigr),
  \qquad
  \alpha = \theta_{\mathrm{cmd}} - \gamma,
  \label{eq:cpr}
\end{equation}
the clamp preventing pitching below the local horizontal. Under the
\texttt{pso\_coast} architecture the pitch rate $\dot\theta$ is exposed as an
additional PSO decision variable rather than being derived from $t_{go}$.

\subsection{Linear Tangent Steering}
\label{ssec:lts}

Linear tangent steering prescribes that the tangent of the velocity inclination
angle $(\alpha+\gamma)$ vary linearly with time-to-go,
$\tan(\alpha+\gamma)=a\,t_{go}+b$. Terminal horizontal flight at $t_{go}=0$ gives
$b=0$, and continuity with the current state gives $a=\tan\gamma/t_{go}$, so that
\begin{equation}
  \alpha = \arctan\!\bigl(a\,t_{go}\bigr) - \gamma,
  \label{eq:lts_alpha}
\end{equation}
recomputed at the guidance update rate. This provides smooth, monotonic
inclination control.

\subsection{Bilinear Tangent Steering}
\label{ssec:bts}

The bilinear law generalises the linear form by writing the inclination tangent
as a ratio of two linear functions of $t_{go}$,
\begin{equation}
  \tan(\alpha+\gamma)
    = \frac{c_1\,t_{go} + c_2}{c_1'\,t_{go} + c_2'},
  \label{eq:bts}
\end{equation}
whose four constants are fixed by four conditions: terminal horizontal flight
($c_2=0$); a small terminal-rate constant $c_1$ for a smooth approach;
normalisation ($c_2'=1$); and an initial-value match to a blended target
$\gamma_{\mathrm{blended}}=0.7\,\gamma_{\mathrm{current}}+0.3\,\gamma_{\mathrm{target}}$,
where $\gamma_{\mathrm{target}}$ is the geometric flight-path angle required to
reach the target altitude over the remaining downrange. Controlling both the
value and the rate of the inclination angle at the terminal point yields a
smoother insertion than the purely linear variant.

\subsection{Apollo Polynomial Guidance}
\label{ssec:apollo}

The Apollo mode commands the total acceleration in the downrange ($x$) and
altitude ($y$) directions with two linear-in-time channels,
$a_x(t)=k_1(t-t_{\mathrm{epoch}})+k_2$ and
$a_y(t)=k_3(t-t_{\mathrm{epoch}})+k_4$, whose coefficients are chosen to reach
the target radius, the local circular velocity $\sqrt{\mu/r_f}$ in the horizontal,
and zero radial velocity at burnout. For a generic channel,
\begin{equation}
  k_1 = \frac{6(v_f+v_i)\,t_{go} - 12(x_f-x_i)}{t_{go}^{3}}, \qquad
  k_2 = \frac{-2(v_f+2v_i)\,t_{go} + 6(x_f-x_i)}{t_{go}^{2}}.
  \label{eq:apollo_k}
\end{equation}
As $t_{go}\to0$ the denominators drive the coefficients to large values, so they
are frozen at their last values once $t_{go}$ falls below a threshold (nominally
$10$~s). The thrust-only acceleration is obtained by subtracting the local gravity
vector, and the steering angle follows as
$\alpha=\arctan(a_y^{\mathrm{thrust}}/a_x^{\mathrm{thrust}})-\arctan(v_y/v_x)$. An
optional thrust-magnitude-control mode throttles the engine to the commanded
acceleration magnitude, up to the maximum available thrust.

\subsection{Powered Explicit Guidance (PEG)}
\label{ssec:peg}

The classical PEG mode maintains the linear pitch program
$\sin(\text{pitch}[t]) = A + B\,t$ and updates the steering constants $A,B$ and the
burn-time estimate $T$ each major cycle. Given the rocket-equation thrust
integrals $b_0,b_1,c_0,c_1$ over the remaining burn, $A$ and $B$ are obtained by
solving the $2\times2$ linear system that enforces the target radius and radial
velocity at burnout, while $T$ is refined from the angular-momentum gap
$h_T-h$~\cite{burden2015numerical}. The undamped Guide$\rightarrow$Estimate
fixed-point iteration can exhibit a two-point limit cycle for early Stage-2
states, so successive under-relaxation (SUR) is applied,
\begin{equation}
  T_{n+1} = \omega\,T_{\mathrm{est}}(T_n) + (1-\omega)\,T_n,
  \qquad \omega=0.5,
  \label{eq:peg_sur}
\end{equation}
which breaks the cycle and converges in a few iterations~\cite{burden2015numerical}.
The minor loop then commands $\alpha=\arcsin(A+B\,\tau)-\gamma$, with $\tau$ the
time since the last major update.

\subsection{Analytical Predictor--Corrector PEG (\texttt{peg\_new})}
\label{ssec:peg_new}

The \texttt{peg\_new} mode implements the velocity-to-be-gained formulation, in
which the primary variable is the two-dimensional vector $\mathbf{v}_{go}$ and the
time-to-go follows directly from the rocket equation,
$t_{go}=\tau\bigl(1-e^{-\lVert\mathbf{v}_{go}\rVert/v_e}\bigr)$. The tangential
component is the (gravity-free) velocity deficit
$v_{go,\theta}=v_{\theta,T}-v_\theta$, while the radial component is iterated
against the gravity integral. A predictor--corrector outer loop refines that
integral by evaluating the radial gravity at the predicted burnout position and
averaging it with the current value (trapezoidal quadrature), which supplies the
physically correct pitch-down for orbit insertion. The steering direction is the
normalised
\begin{equation}
  \hat{\mathbf u}(t) \;\propto\; \frac{\mathbf{v}_{go}}{L_0}
        + \lambda'_r\,(t-t_\lambda)\,\hat{\mathbf r},
  \label{eq:pegnew_u}
\end{equation}
with $L_0=\lVert\mathbf{v}_{go}\rVert$, the radial position costate $\lambda'_r$,
and the reference time $t_\lambda$ obtained from the constant-thrust integrals;
the angle of attack is $\alpha=\beta-\gamma$ with
$\beta=\operatorname{atan2}(u_r,u_\theta)$. This gravity-aware $t_{go}$ is also
reusable as the \texttt{TGO\_ESTIMATOR} for the scalar-$t_{go}$ laws.

\subsection{Exponential Pitch (Shooting) Law}
\label{ssec:exp_shooting}

The exponential pitch law prescribes an open-loop pitch program
$\theta(t_{\mathrm{rel}}) = a\,e^{b\,t_{\mathrm{rel}}}$, so that
$\alpha = a\,e^{b\,t_{\mathrm{rel}}}-\gamma$. The pair $(a,b)$ is fixed once, at
guidance start, by a single-shooting solve: \texttt{scipy.optimize.fsolve} adjusts
$(a,b)$ until a simplified drag-free forward propagation to burnout satisfies the
two terminal constraints $r(T)=r_T$ and $\gamma(T)=0$. If the solve fails to
converge the current-state initial guess is retained. This is the simplest
boundary-value method, and the same shooting idea underlies the indirect
optimal-control architecture, where the unknowns are the initial costates rather
than pitch coefficients.

\subsection{Indirect Optimal-Control Steering (\texttt{indirect\_pmp})}
\label{ssec:indirect_guidance}

The indirect mode commands the costate-optimal angle of attack derived in
Section~\ref{ssec:indirect_ascent},
$\alpha=\operatorname{atan2}\!\bigl(-\lambda_\gamma/v,\,-\lambda_v\bigr)$
(Eq.~\eqref{eq:indirect_control}), obtained by propagating the augmented
state--costate system alongside the trajectory. Unlike the other eight modes it is
not a stand-alone feedback law: the initial costates that make it an extremal are
themselves found by the dedicated optimisation architecture of
Section~\ref{sssec:indirect_pmp}, which also governs how the law is replayed when
it steers a segment of a multi-law schedule.

% =============================================================================
\section{Mission Architectures and Optimization Strategies}
\label{sec:optimization}

The optimiser is deliberately guidance-agnostic: it supplies a design vector and
receives a scalar cost, while the inner simulation selects and runs whichever
guidance law is active. Four architectures are available, selected by
\texttt{COAST\_METHOD} together with the \texttt{indirect\_pmp} mode. All share
the non-dimensional penalised-objective philosophy introduced in
Section~\ref{ssec:optimization}: terminal errors and infeasibility are folded
into a single scalar to be minimised.

\subsection{\texttt{apogee\_check}: Brute-Force Kick-Angle Search}
\label{ssec:apogee_check}

The \texttt{apogee\_check} architecture optimises the single kick angle $\alpha$
explicitly while determining the cut-off implicitly through a physics-based event.
The objective is the total propellant required to reach the target circular orbit,
\begin{equation}
  \min_{\alpha}\; J(\alpha), \qquad
  J(\alpha)=m_{\mathrm{prop,ascent}}(\alpha)+m_{\mathrm{prop,circ}}(\alpha),
  \label{eq:apogee_objective}
\end{equation}
subject to reaching the target apogee radius within the propellant and
burn-duration limits. The circularisation propellant is evaluated from the
Tsiolkovsky equation,
\begin{equation}
  m_{\mathrm{prop,circ}} = m_{\mathrm{apo}}
    \Bigl(1 - \exp\!\bigl[-\Delta v_{\mathrm{circ}}/(I_{sp}g_0)\bigr]\Bigr),
  \qquad
  \Delta v_{\mathrm{circ}} = \bigl|v_{\mathrm{circular}}-v_{\mathrm{apo}}\bigr|,
  \label{eq:circ_dv}
\end{equation}
with $v_{\mathrm{circular}}=\sqrt{\mu/r_{\mathrm{target}}}$. Above the
atmosphere-exit altitude the powered phase is terminated when the
\emph{instantaneous unpowered} apogee of the current state matches the target,
detected through the event function
$g(t,\mathbf{y}) = r_{\mathrm{apo}}(t,\mathbf{y}) - r_{\mathrm{target}}$, where
$r_{\mathrm{apo}}=a(1+e)$ is computed from the instantaneous orbital elements
$\varepsilon=v^{2}/2-\mu/r$, $a=-\mu/2\varepsilon$,
$e=\sqrt{1-h^{2}/(\mu a)}$~\cite{curtis2020orbital}. Because the design space is
one-dimensional, the optimum $\alpha^{\star}=\arg\min_{\alpha_i}J(\alpha_i)$ is
found by a brute-force grid search (\texttt{scipy.optimize.brute}), which is
robust to the non-smoothness introduced by discrete staging and cut-off events.

\subsection{\texttt{pso\_coast}: Coast-Parameter PSO}
\label{ssec:pso_coast}

The \texttt{pso\_coast} architecture searches a four-variable design vector with
particle swarm optimisation~\cite{kennedy1995particle,biscani2020pagmo},
\begin{equation}
  \mathbf{x} = \bigl[\,\Delta t_c,\; \Delta t_{r}\%,\; \text{coast-start}\%,\;
                     \gamma_p\,\bigr],
  \label{eq:pso_coast_vector}
\end{equation}
where $\Delta t_c$ is the coast duration, $\Delta t_r\%$ the Stage-2 burn as a
percentage of the maximum burn time, coast-start$\%$ the coast onset as a
percentage of that burn, and $\gamma_p$ the pitch-manoeuvre angle
(kick angle $=\gamma_p-\pi/2$). The trajectory is assembled as a
thrust--coast--thrust sequence of three arcs sharing the state, and scored by a
four-term non-dimensional objective (a burn-time term plus the normalised
altitude, velocity, and flight-path-angle errors) augmented by a large fixed
crash penalty --- the same normalisation as the indirect objective of
Section~\ref{sssec:indirect_pmp} without its transversality term.

\subsection{\texttt{direct}: Direct-Insertion PSO}
\label{ssec:direct}

The \texttt{direct} architecture inserts the payload with a single continuous
Stage-2 burn cut at circular velocity, with no coast-to-apogee and no
circularisation burn. It is a two-variable PSO over
\begin{equation}
  \mathbf{x} = \bigl[\,\gamma_p,\; t_{\mathrm{burn}}\%\,\bigr],
  \label{eq:direct_vector}
\end{equation}
the pitch angle and the Stage-2 burn duration as a percentage of the maximum. The
objective is the same non-dimensional weighted sum,
\begin{equation}
  J = w_J\,\frac{t_{\mathrm{burn}}}{T_{\max,2}}
    + w_h\,\bigl|\Delta h\bigr|_{\mathrm{nd}}
    + w_v\,\bigl|\Delta v\bigr|_{\mathrm{nd}}
    + w_\gamma\,\bigl|\Delta\gamma\bigr|_{\mathrm{nd}} + C,
  \label{eq:direct_objective}
\end{equation}
where the errors are measured against a direct-insertion box target and $C$ is the
crash penalty. Only a subset of the guidance laws (notably Apollo, PEG, and
\texttt{peg\_new}) can reliably close a single continuous burn; others go
suborbital under this architecture.

\subsection{\texttt{indirect\_pmp}: Indirect Optimal-Control PSO}
\label{sssec:indirect_pmp}

The \texttt{indirect\_pmp} architecture realises the optimal-control reference
trajectory formulated in Section~\ref{ssec:indirect_ascent}. It propagates the
augmented state--costate system --- the reduced state $[s,r,v,\gamma,m]$ together
with the adjoints $[\lambda_r,\lambda_v,\lambda_\gamma]$ governed by
Eqs.~\eqref{eq:indirect_costates} and the costate-optimal control
$\alpha=\operatorname{atan2}(-\lambda_\gamma/v,-\lambda_v)$ of
Eq.~\eqref{eq:indirect_control} --- and searches the unknown initial costates and
arc-timing parameters with a particle swarm optimiser.

\paragraph{Decision vector.}
Each particle encodes the seven-element vector
\begin{equation}
  \mathbf{x} = \bigl[\,\lambda_{0r},\;\lambda_{0v},\;\lambda_{0\gamma},\;
        \Delta t_c,\; \text{coast-start},\; \Delta t_T,\; \gamma_p\,\bigr],
  \label{eq:indirect_vector}
\end{equation}
the initial costates, the coast duration, the coast-start (initialisation) time,
the last-stage burn duration, and the pitch-manoeuvre angle. Because the costate
equations are homogeneous in $\boldsymbol\lambda$, only the \emph{direction} of
the initial costate matters, so the three costate components are searched on the
normalised interval $[-1,1]$ and the $\lVert\boldsymbol\lambda\rVert=1$ gauge is
fixed at the pitch-over point where the transversality residual is evaluated.

\paragraph{Augmented objective.}
The boundary and optimality conditions are enforced through a static penalty,
turning the two-point boundary-value problem into an unconstrained minimisation,
\begin{equation}
  J' = J + \sum_{c=1}^{3} s_c\,\bigl\lVert x_{c,f}-x_c'\bigr\rVert
     + s_4\,\bigl\lVert H_f^{\mathrm{last}} + H_f^{\mathrm{coast}}
                      - H_0^{\mathrm{last}}\bigr\rVert + C,
  \label{eq:indirect_augmented}
\end{equation}
where $J=t_f-t_{cf}$ is the powered-flight time, the first sum penalises the
terminal altitude, velocity, and flight-path-angle errors, the fourth term
penalises violation of the transversality condition
Eq.~\eqref{eq:indirect_transversality}, and $C$ is a large fixed penalty
($10^{20}$) applied to non-physical (crashing or propellant-depleting)
trajectories. Minimising $J'$ drives the swarm towards feasible extremals; the
run uses a swarm of $250$ particles over $500$ generations.

\paragraph{Multi-arc integration.}
When the optional full-ascent extension is active, the trajectory is integrated as
one continuous pass of the augmented state through an ordered arc list --- a
vertical rise at $\alpha\equiv0$, the first-stage powered arc terminated by a
propellant-depletion (MECO) event, an instantaneous inert-mass staging drop, an
inter-stage coast, and the upper-stage thrust--coast--thrust arcs. The costates are
seeded at the pitch-over point (the homogeneous adjoint equations keep them zero
through the vertical rise), and, since the reduced costate set carries no mass
adjoint, the position and velocity costates and the Hamiltonian remain continuous
across the staging and coast corners by the Weierstrass--Erdmann conditions, so no
interior jump conditions are imposed. The angle-of-attack clamp of
Eq.~\eqref{eq:fa_alpha_clamp} is active only while the vehicle is inside the
atmosphere, using the same \texttt{ATMOSPHERE\_EXIT\_METHOD} criterion as the
fairing-jettison and guidance-activation logic, and is lifted at atmosphere exit so
the near-vacuum arc recovers the exact interior control. The passive mass costate
$\lambda_m$ of Eq.~\eqref{eq:fa_masscostate} is integrated for rigour; because it
feeds back into nothing, its terminal condition $\lambda_m(t_f)=0$ is imposed in
closed form by subtracting the constant $\lambda_m(t_f)$ from the integrated
history, which renders the Hamiltonian constant along each powered arc.

\subsection{Separation of Concerns}
\label{ssec:separation}

Across all four architectures the outer optimiser supplies only a design vector
and receives only a scalar cost, while the inner simulation selects and runs
whichever guidance law is active. This decoupling is what makes the comparative
study fair: the same trajectory engine, environment models, and scoring
normalisation are applied to every guidance law, so differences in the results
reflect the laws themselves rather than the machinery around them.

% =============================================================================
\section{Configuration System}
\label{sec:configuration}

All user-facing behaviour is controlled from
\texttt{Input\_File/simulation\_parameters.py}, without editing the solver code.
Table~\ref{tab:config} groups the principal switches by function; each mission is
fully specified by a choice within these groups plus the vehicle and
mission-geometry parameters of Section~\ref{sec:vehicle}.

\begin{table}[htbp]
  \centering
  \caption{Principal configuration switches, grouped by function.}
  \label{tab:config}
  \small
  \begin{tabular}{l l}
    \toprule
    \textbf{Group} & \textbf{Representative parameters} \\
    \midrule
    Guidance selection      & \texttt{GUIDANCE\_MODE} (nine laws),
                              \texttt{TGO\_ESTIMATOR}, \texttt{PEG\_CONVERGENCE\_MODE},
                              \texttt{CPR\_THETA\_DOT\_MODE} \\
    Mission architecture    & \texttt{COAST\_METHOD}
                              (\texttt{apogee\_check}/\texttt{pso\_coast}/\texttt{direct}) \\
    Earth-rotation / frames & \texttt{ENABLE\_EARTH\_ROTATION},
                              \texttt{INCLUDE\_PSEUDO\_FORCES},
                              \texttt{AZIMUTH\_INCLINATION\_MODE} \\
    Atmosphere / aero       & \texttt{INCLUDE\_DRAG}, \texttt{INCLUDE\_LIFT},
                              \texttt{ATMOSPHERE\_EXIT\_METHOD} \\
    Kick profile            & \texttt{KICK\_PROFILE\_MODE} (triangular / instantaneous) \\
    Engine modes            & \texttt{ISP\_1\_MODE}, \texttt{THRUST\_1\_MODE}
                              (sea-level / vacuum / average / linear) \\
    Optimiser settings      & \texttt{RUN\_FAST}, PSO tuning blocks
                              ($\omega$, $c_1$, $c_2$, $v_{\max}$, swarm size,
                              generations, seed) \\
    Integration / output    & tolerances, maximum step, dense-output resolution,
                              plot flags \\
    \bottomrule
  \end{tabular}
\end{table}

% =============================================================================
\section{Vehicle Configuration and Baseline Mission}
\label{sec:vehicle}

The simulator is configured with parameters representative of a Falcon-9-class
two-stage launcher (Table~\ref{tab:stage_params}). The framework is not specific
to this vehicle: any two-stage launcher can be studied by editing the stage
masses, engine parameters, reference area, and drag coefficient; the Falcon-9
values are used here as a concrete, well-documented baseline.

\begin{table}[ht]
\centering
\caption{Launch vehicle stage parameters~\cite{NASA_Falcon9v12_Datasheet_2018,
SpaceX_Falcon9_2025, Wevolver_Falcon9v12_Block5_2019}.}
\label{tab:stage_params}
\begin{tabular}{lcc}
\hline
\textbf{Parameter} & \textbf{First Stage} & \textbf{Second Stage} \\
\hline
Isp (Sea Level)          & 283 s        & --       \\
Isp (Vacuum)             & 312 s        & 348 s    \\
Maximum Thrust (SL)      & 7{,}607 kN    & 8{,}227 kN \\
Maximum Thrust (V)       & --           & 981 kN   \\
Propellant               & LOX/RP-1      & LOX/RP-1 \\
Propellant Mass          & 395{,}700 kg  & 92{,}670 kg \\
Empty Mass               & 25{,}600 kg   & 3{,}900 kg \\
Payload Capacity (LEO)   & --           & 22{,}800 kg \\
Payload Capacity (GTO)   & --           & 8{,}300 kg \\
Diameter                 & 3.7 m        & 3.7 m    \\
\hline
\end{tabular}
\end{table}

The launch site and target orbit are set independently by the user
(Table~\ref{tab:param_mission}); the values below are the baseline used for the
demonstration case of the next chapter. The baseline targets a $500$~km circular
orbit at an inclination of $51.6^\circ$ (representative of the International Space
Station), launched from latitude $28.5^\circ$ (Kennedy Space Center). The
resulting lift-off mass is approximately $517.9$~t, being the sum of both stages
with zero payload in the reference configuration.

\begin{table}[htbp]
  \centering
  \caption{Mission geometry and target-orbit parameters.}
  \label{tab:param_mission}
  \begin{tabular}{l l c c}
    \toprule
    \textbf{Parameter} & \textbf{Description} & \textbf{Default} & \textbf{Unit} \\
    \midrule
    Target orbital altitude    & Circular orbit height above surface
                               & 500       & km   \\
    Target inclination         & Desired final orbital inclination
                               & 51.6      & deg  \\
    Launch latitude            & Geodetic latitude of the launch pad
                               & 28.5      & deg  \\
    Launch longitude           & Longitude of the launch pad (reserved)
                               & $-80.5$   & deg  \\
    \bottomrule
  \end{tabular}
\end{table}

% =============================================================================
\section{Software Usage}
\label{sec:software}

A single mission is configured entirely in
\texttt{Input\_File/simulation\_parameters.py} and run through \texttt{main.py},
which selects the guidance law and mission architecture, runs the appropriate
optimiser, propagates the optimal trajectory, and writes the result arrays and
summary metrics. Vehicle and environment models are taken from the
\texttt{Auxiliary/} package, and no source-code edits are needed to change the
vehicle, mission, guidance law, or optimiser.

For comparative studies, batch runners execute every guidance law under common
conditions and collate the outcomes
(\texttt{all\_guidance\_plotting/run\_all\_guidance\_methods.py} and
\texttt{guidance\_comparison/compare\_guidance\_methods.py}). The plotting layer
(\texttt{Plots/new\_plot\_runner.py}, drawing on \texttt{Plots/new\_metrics/})
consumes the plain result arrays produced by any run and is independent of the
solvers, so figures can be regenerated without re-running the simulation.
